\emph{Charge-to-edge moment identification in the Gaussian rank-one class.} Within the Gaussian state model of the Row length paragraph - independent momenta with a variance profile, a cap block with block-row profile \(g\) of total weight \(W\) and edge datum \(s^{2}\), and the leading-order convention neglecting \(\operatorname{Im}\alpha\) - the identification of the moments of the classicalised edge charge with the edge data of the block cut, which the transplanted selection rule uses as input, is an exact moment map with one quantified profile obstruction. At the leading-order convention the nonzero eigenvalues of a single cap's charge are exactly \(\pm\ell\) with \(\ell=2|w_\perp|\) the row length, so the charge moment sequence is exactly \(2\ell^{2p}\) and all odd charge moments vanish, pointwise in the configuration; for the two band-separated caps of one axis the spectrum is the union of the two pairs and the moment sequences add. Over the Gaussian state, the averaged second moment is \(2(4s^{2}W)=8s^{2}W\), consistent through \(W\) alone with \(\langle\ell^{2}\rangle=4s^{2}W\) of the Row length paragraph, so the identification at the second moment - the one that fixes the row length - is exact in the model and carries no state dependence beyond the row-sum \(W\). Beyond second order the identification is obstructed exactly and strictly: with \(D_{\mathrm{eff}}=(\sum g)^{2}/\sum g^{2}\) and \(\rho_{p}(g)=p!\,h_{p}(g)/(\sum g)^{p}\), the averaged \(2p\)-th moment equals the edge-data single-scale prediction multiplied by \(\rho_{p}(g)\), and \(\rho_{p}-1\) is bounded below by \(C(p,2)/D_{\mathrm{eff}}\), equals \(1/D_{\mathrm{eff}}\) at \(p=2\), equals \(3/D_{\mathrm{eff}}+2s_{3}/(\sum g)^{3}\) at \(p=3\) with \(s_{3}=\sum_i g_i^{3}\), obeys \(1<\rho_{p}\le p!\) with the upper end attained at a single row, and has the exact partition form \(\rho_{p}=\sum_{\lambda\vdash p}(p!/z_{\lambda})\prod_{k\ge1}\bigl(s_{k}/(\sum g)^{k}\bigr)^{m_{k}}\) in terms of the power sums \(s_{k}\) and the multiplicities \(m_{k}\) of the parts of \(\lambda\vdash p\). The identification from edge data alone is therefore exact at the second moment, fails at the fourth by exactly \(1/D_{\mathrm{eff}}\), and is recovered at leading order only as \(D_{\mathrm{eff}}\to\infty\): for the fuzzy variance proxies \(f=\nu l/2\) the effective dimension grows with the cutoff, from \(1488\) at \(l_c=40\) to \(23017\) at \(l_c=160\), while the fourth-moment excess drops from \(4.0\times10^{-3}\) to \(2.6\times10^{-4}\) (numerical, on the exact rational profiles). The per-factor moment-signature identity of the mixed-multiplicity paragraph is the zero-imaginary-part case of the same map, reproduced from full charge spectra with excesses \(9/17\) and \(81/145\) at \((q_{1},q_{2})=(2,1)\), \((n_{1},n_{2})=(1/4,\,2/9)\). These statements are for the rank-one edge-datum family with independent Gaussian momenta and block-row profiles, at the leading-order convention neglecting \(\operatorname{Im}\alpha\), with band-separated caps and fixed multiplicities, inside the transplanted rule with its single-dominant-diagram hypothesis at leading semiclassical order and its diagonal-cut convention; the obstruction factor is exact for every finite profile, the identification from edge data holds exactly only at the second moment and at leading order in the concentrated-profile regime, they do not fix the state that sets \(W\) and \(g\), they leave the selection rule, its transplant and the missing steps (i)-(iv) of the Missing step paragraph unchanged, and they close nothing in Problem~G.
